% Generated by tools/build_new_dists_anthology_sections.py; do not edit by hand.

\section{Central extensions and projective representations}\label{proof:new-dist:central_extensions}

\resultbox{\textbf{Canonical testing artifacts.}\par\smallskip Exact backend: \path{lambda_distributions/dists/central_extensions_projective_representations/verification.py}.\par Regression: \path{lambda_distributions/dists/central_extensions_projective_representations/test_verification.py}.\par Marimo: \path{marimo_notebooks/central_extensions.py}.}

\subsection*{Scope and outcome}

For a finite or compact central extension acting with one scalar central
character, we prove that the one-sided generalized Molien series is obtained
by a homogeneous root-of-unity projection.  A continuous scalar center kills
every positive degree.  Since a projective matrix has no canonical
one-sided eigenvalue distribution, we construct and prove lift independence
of the charge-zero balanced series.  We then specialize the theory to Schur
and pin covers, finite Weil representations, extraspecial groups, Clifford
groups, sporadic covers, and determinant-one kernels.  The companion marimo
notebook constructs nine representative matrix groups in dimensions $2$,
$3$, $5$, and $9$.  It performs 180 coefficient comparisons through total
degree six, five full determinant-average checks, and nine balanced-moment
checks.  Every comparison passes; the maximum numerical discrepancy is
$6.28\times10^{-12}$.

\subsection{The generalized Molien target}

Let $K$ be a finite or compact group and let
$\rho:K\to U(V)$ be a finite-dimensional unitary representation.  For a
partition $\tau=(\tau_1,\ldots,\tau_\ell)$ set
\[
 W_\tau(V)=\bigotimes_{i=1}^{\ell}\Sym^{\tau_i}V,
 \qquad |\tau|=\sum_i\tau_i.
\]
The one-sided $\sigma$-MGF is
\begin{equation}\label{nd:central_extensions:eq:molien}
 \MGF_{K,V}(\mathbf t)
 =\int_K\prod_{i\geq1}\det(I-t_i\rho(g))^{-1}\,dg.
\end{equation}

\begin{theorem}[Generalized Molien identity]\label{nd:central_extensions:thm:molien}
With normalized counting measure in the finite case and normalized Haar
measure in the compact case,
\[
 \MGF_{K,V}(\mathbf t)
 =\sum_\tau\dim W_\tau(V)^K\,m_\tau(\mathbf t).
\]
Equivalently,
\begin{equation}\label{nd:central_extensions:eq:trace-exp}
 \prod_i\det(I-t_i\rho(g))^{-1}
 =\exp\!\left(\sum_{k\geq1}\frac{p_k(\mathbf t)}{k}
                   \Tr(\rho(g)^k)\right).
\end{equation}
\end{theorem}

\begin{proof}
For every matrix $A$,
$\det(I-tA)^{-1}=\sum_{d\geq0}\Tr(\Sym^d A)t^d$.
Multiplication in the variables $t_i$ makes the coefficient of
$m_\tau$ the character of $W_\tau(V)$.  Haar averaging is the Reynolds
projection, whose trace is the dimension of the fixed space.  Finally,
$-\log\det(I-tA)=\sum_{k\geq1}t^k\Tr(A^k)/k$ proves
\eqref{nd:central_extensions:eq:trace-exp}.
\end{proof}

\subsection{Central-character projection}

Consider a central extension
\begin{equation}\label{nd:central_extensions:eq:extension}
 1\longrightarrow Z\longrightarrow\widetilde G
 \xrightarrow{\pi}G\longrightarrow1
\end{equation}
and a representation $\rho:\widetilde G\to U(V)$ for which
$\rho(z)=\chi(z)I_V$.  Write $C=\chi(Z)\leq U(1)$.

\begin{theorem}[Finite central-character projector]\label{nd:central_extensions:thm:projector}
Suppose $C=\mu_r$.  Then
\begin{equation}\label{nd:central_extensions:eq:selection}
 [m_\tau]\MGF_{\widetilde G,V}=0\qquad\text{unless }r\mid|\tau|.
\end{equation}
For any set-theoretic section $g\mapsto\widetilde g$,
\begin{equation}\label{nd:central_extensions:eq:fibre-average}
 \boxed{\MGF_{\widetilde G,V}(\mathbf t)
 =\frac1{|G|}\sum_{g\in G}\frac1r\sum_{\zeta\in\mu_r}
 \prod_i\det(I-\zeta t_i\rho(\widetilde g))^{-1}.}
\end{equation}
Thus adjoining scalar $r$th roots applies
\begin{equation}\label{nd:central_extensions:eq:Pr}
 \PP_rF(\mathbf t)=\frac1r\sum_{\zeta\in\mu_r}F(\zeta\mathbf t),
\end{equation}
which retains precisely the homogeneous degrees divisible by $r$.
\end{theorem}

\begin{proof}
The center acts on $W_\tau(V)$ by the scalar $\chi(z)^{|\tau|}$.
A fixed vector can exist only when this scalar is one for all $z$, proving
\eqref{nd:central_extensions:eq:selection}.  For the integral proof, disintegrate normalized
measure over the fibres of $\pi$.  The pushforward of normalized measure on
$Z$ under $\chi$ is uniform measure on $C$.  In total degree $d$, the fibre
average is
\[
 \frac1r\sum_{\zeta\in\mu_r}\zeta^d
 =\begin{cases}1,&r\mid d,\\0,&r\nmid d.\end{cases}
\]
This proves \eqref{nd:central_extensions:eq:fibre-average} and \eqref{nd:central_extensions:eq:Pr}, independently of
the chosen section and even when $\chi$ has a kernel.
\end{proof}

\begin{corollary}[Continuous scalar saturation]\label{nd:central_extensions:cor:u1}
If $C=U(1)$, then
\[
 \boxed{\MGF_{\widetilde G,V}=1.}
\]
\end{corollary}

\begin{proof}
The scalar average in total degree $d$ is
$\int_{U(1)}z^d\,dz$, equal to zero for every $d>0$ and one for $d=0$.
\end{proof}

This explains the trivial one-sided series of $U(d)$ in its defining
module and of a full unitary normalizer containing every scalar matrix.

\subsection{The canonical projective series}

A projective matrix $[U_g]\in PU(V)$ has no canonical eigenvalue multiset:
another lift $\zeta_gU_g$ multiplies every eigenvalue by $\zeta_g$.
Consequently a projective representation does not canonically define
\eqref{nd:central_extensions:eq:molien}.

For partitions $\alpha,\beta$, define the charge-zero balanced series
\begin{equation}\label{nd:central_extensions:eq:balanced-def}
 \boxed{\BSMG^{\mathrm{proj}}_{G,[V]}(\mathbf s,\mathbf t)
 =\sum_{|\alpha|=|\beta|}
 \dim\!\left(W_\alpha(V)\otimes W_\beta(V)^*\right)^{\widetilde G}
 m_\alpha(\mathbf s)m_\beta(\mathbf t).}
\end{equation}

\begin{theorem}[Lift-independent projective Molien formula]
For any choice of unitary lift $U_g$ of each projective matrix,
\begin{equation}\label{nd:central_extensions:eq:balanced-ct}
 \boxed{\BSMG^{\mathrm{proj}}_{G,[V]}
 =\CT_u\int_G
 \prod_i\det(I-u s_iU_g)^{-1}
 \prod_j\det(I-u^{-1}t_jU_g^{-*})^{-1}\,dg.}
\end{equation}
The answer is unchanged by arbitrary, element-dependent replacements
$U_g\mapsto\zeta_gU_g$.  It therefore depends only on the projective image,
not on a Schur cover, cocycle representative, scalar phases, or a
determinant-one lift.
\end{theorem}

\begin{proof}
The coefficient indexed by $(\alpha,\beta)$ in the integrand acquires
$u^{|\alpha|-|\beta|}$.  The constant term imposes equality of the two
degrees.  Under $U_g\mapsto\zeta_gU_g$, that coefficient is multiplied by
$\zeta_g^{|\alpha|-|\beta|}=1$.  Haar averaging then gives the invariant
dimension exactly as in Theorem~\ref{nd:central_extensions:thm:molien}.
\end{proof}

For a finite scalar center $\mu_r$, the full two-sided series has sectors
$|\alpha|-|\beta|\equiv0\pmod r$.  Exact equality, rather than congruence,
is the part intrinsic to the projective representation.  The multilinear
diagonal coefficients are the frame potentials
\begin{equation}\label{nd:central_extensions:eq:frame}
 [m_{(1^k)}(\mathbf s)m_{(1^k)}(\mathbf t)]\BSMG^{\mathrm{proj}}
 =\int_G|\Tr U_g|^{2k}\,dg.
\end{equation}

\subsection{Finite groups: class compression and power maps}

If $\widetilde G$ is finite with character $\chi_V$, grouping
\eqref{nd:central_extensions:eq:trace-exp} by conjugacy classes gives
\begin{equation}\label{nd:central_extensions:eq:class-formula}
 \boxed{\MGF_{\widetilde G,V}
 =\frac1{|\widetilde G|}\sum_{C}|C|
 \exp\!\left(\sum_{k\geq1}\frac{p_k(\mathbf t)}{k}
                         \chi_V(c_C^k)\right).}
\end{equation}
Thus an exact calculation needs the represented character, class sizes, and
class power maps.  The projective balanced version is
\begin{equation}\label{nd:central_extensions:eq:balanced-class}
 \CT_u\frac1{|\widetilde G|}\sum_C|C|\exp\!\left[
 \sum_{k\geq1}\frac{u^kp_k(\mathbf s)\chi_V(c_C^k)
 +u^{-k}p_k(\mathbf t)\overline{\chi_V(c_C^k)}}{k}\right].
\end{equation}
Equations \eqref{nd:central_extensions:eq:class-formula} and \eqref{nd:central_extensions:eq:balanced-class} are the
practical route for large Schur covers, spin Weyl covers, and sporadic
covers whose dense matrices would be wasteful.

\subsection{Schur and spin covers: the binary octahedral test}

For a genuine spin module of a double cover $2.S_n$ or $2.A_n$, the central
involution acts by $-I$.  Theorem~\ref{nd:central_extensions:thm:projector} immediately gives
\begin{equation}\label{nd:central_extensions:eq:spin-parity}
 [m_\tau]\MGF_{2.S_n,V_n}=[m_\tau]\MGF_{2.A_n,V_n}=0
 \qquad(|\tau|\text{ odd}).
\end{equation}

For the standard basic-spin module, Clifford algebra identifies
$\End(V_n)$ with all or one parity of the exterior algebra of the reflection
module $E_n$.  Its exterior powers are distinct hook modules.  With the
standard associate convention this gives
\begin{equation}\label{nd:central_extensions:eq:basic-spin-fourth}
 \frac1{|2.S_n|}\sum_g|\chi_{V_n}(g)|^4
 =\begin{cases}n,&n\text{ odd},\\n/2,&n\text{ even}.
 \end{cases}
\end{equation}
The fourth balanced moment therefore grows linearly although $\dim V_n$
grows exponentially; it does not approach the Haar value $2$.

The notebook uses the first nontrivial concrete model relevant here:
the $48$ unit quaternions of the binary octahedral group.  Under
\[
 a+bi+cj+dk\longmapsto
 \begin{pmatrix}a+bi&c+di\\-c+di&a-bi\end{pmatrix}\in SU(2),
\]
they form a double cover $2.S_4\to S_4$ and give the two-dimensional
basic-spin representation.  The list consists of
\[
 \{\pm1,\pm i,\pm j,\pm k\},\qquad
 \tfrac12(\pm1\pm i\pm j\pm k),\qquad
 \tfrac1{\sqrt2}(\pm e_a\pm e_b)\ (a<b).
\]
All signs are independent in the displayed families.  Dense multiplication
checks all $48^2$ products.  The fourth trace moment is $2$, agreeing with
\eqref{nd:central_extensions:eq:basic-spin-fourth} at $n=4$.

\begin{center}
\begin{tabular}{@{}c c c c c@{}}
\toprule
$k$ & direct $|\Tr|^{2k}$ & phase-twisted lifts & Haar $U(2)$ & conclusion\\
\midrule
1 & 1 & 1 & 1 & agree\\
2 & 2 & 2 & 2 & agree\\
3 & 5 & 5 & 5 & agree\\
4 & 15 & 15 & 14 & first separation\\
\bottomrule
\end{tabular}
\end{center}

The phase-twisted column multiplies the $j$th matrix by an unrelated
$37$th root of unity.  This is a direct test that projective balanced
quantities do not depend on a coherent linear lift.  The non-multilinear
test
\[
 (\alpha,\beta)=((2,1),(1,1,1))
\]
also remains $3$ after phase twisting, checking charge-zero invariance beyond
frame potentials.

\subsection{Pin covers of Weyl groups}

Let $W\leq O(E)$ be a finite Weyl or Coxeter group,
$\widetilde W\leq\operatorname{Pin}(E)$ its pin cover, and $S$ a genuine
spinor.  The central involution gives the parity rule
\[
 [m_\tau]\MGF_{\widetilde W,S}=0\qquad(|\tau|\text{ odd}).
\]
Moreover
\begin{equation}\label{nd:central_extensions:eq:cliff-end}
 \End(S)\cong\Cl(E)\ \text{or}\ \Cl^0(E)
 \cong\bigwedge^\bullet E\ \text{or}\ \bigwedge^{\mathrm{even}}E
\end{equation}
as $W$-modules.  Hence
\begin{equation}\label{nd:central_extensions:eq:weyl-moment}
 \boxed{\frac1{|\widetilde W|}\sum_{\widetilde w}
 |\Tr S(\widetilde w)|^{2k}
 =\dim\left(\End(S)^{\otimes k}\right)^W.}
\end{equation}
For $k=2$, if $\End(S)=\bigoplus_\lambda m_\lambda U_\lambda$, the answer
is $\sum_\lambda m_\lambda^2$.  Distinct exterior powers in the classical
towers therefore produce a rank-growing exterior-algebra law, not a Haar
unitary law.  The notebook's $Q_8$ and binary octahedral constructions are
the rank-small pin/spin test cases; higher ranks are more efficiently tested
by \eqref{nd:central_extensions:eq:cliff-end} or class power maps.

\subsection{Finite Weil representations}

Let $W_n=\F_q^{2n}$ and let $\mathcal H_n$ be the full finite
Schr\"odinger space, $\dim\mathcal H_n=q^n$.  For Weyl operators
$\{D_v:v\in W_n\}$, metaplectic covariance says
\[
 U_gD_vU_g^{-1}=D_{gv}
\]
after a compatible normalization.  Since the $D_v$ form a basis of
$\End(\mathcal H_n)$,
\begin{equation}\label{nd:central_extensions:eq:weil-end}
 \End(\mathcal H_n)\cong\C[W_n]
\end{equation}
as an $\Sp(2n,q)$-module.  Burnside's orbit lemma gives
\begin{equation}\label{nd:central_extensions:eq:weil-orbits}
 \boxed{\frac1{|\Sp(2n,q)|}\sum_g|\Tr U_g|^{2k}
 =\#\bigl(\Sp(2n,q)\backslash W_n^k\bigr).}
\end{equation}

There are two orbits on $W_n$: zero and nonzero, so the second trace moment
is $2$.  For ordered pairs and $n\geq2$, the orbit types are
\begin{itemize}
\item $(0,0)$, $(0,w\ne0)$, and $(v\ne0,0)$;
\item $w=av$, one orbit for each $a\in\F_q^\times$;
\item independent orthogonal pairs;
\item independent pairs with $\langle v,w\rangle=c$, one orbit for each
      $c\in\F_q^\times$.
\end{itemize}
Thus the stable fourth moment is $3+(q-1)+1+(q-1)=2q+2$.
Witt extension proves that every fixed-$k$ orbit count stabilizes once the
ambient symplectic rank is large enough.

\begin{remark}[Rank-one correction]
When $n=1$, an orthogonal pair in a two-dimensional symplectic space is
dependent.  The independent-orthogonal orbit is absent, so
\begin{equation}\label{nd:central_extensions:eq:rank-one-weil}
 \mathbb E|\Tr U_g|^4=2q+1\qquad(n=1).
\end{equation}
This hypothesis is essential when testing a reasonable small dimension.
\end{remark}

For odd prime $q$, the notebook independently constructs the projective Weil
matrices from
\[
 F_{xy}=q^{-1/2}e^{2\pi ixy/q},\qquad
 P_{xx}=e^{2\pi i(2^{-1}x^2)/q}.
\]
Projective closure yields $|SL(2,3)|=24$ matrices in dimension $3$ and
$|SL(2,5)|=120$ matrices in dimension $5$.  Separately, the code enumerates
the $SL(2,q)$-orbits on $(\F_q^2)^k$.

\begin{center}
\begin{tabular}{@{}c c c c c@{}}
\toprule
$q$ & projective order & $k$ & matrix average & orbit count\\
\midrule
3 & 24  & 1 & 2 & 2\\
3 & 24  & 2 & 7 & 7\\
5 & 120 & 1 & 2 & 2\\
5 & 120 & 2 & 11 & 11\\
\bottomrule
\end{tabular}
\end{center}

For an irreducible even or odd Weil constituent, one must insert the
corresponding idempotent in \eqref{nd:central_extensions:eq:weil-end}; the full-module orbit count
must not be copied unchanged.  Likewise, over finite fields where a
metaplectic extension splits, the actual scalar image, rather than the name
of the cover, determines the selection rule.

\subsection{Odd extraspecial groups and their central products}

Let $E_n(p)$ be an extraspecial group of order $p^{1+2n}$ and exponent $p$,
where $p$ is odd.  Its Schr\"odinger module has dimension $N=p^n$.  In the
explicit model used by the notebook, for $a,b\in\F_p^n$ and $c\in\F_p$,
\begin{equation}\label{nd:central_extensions:eq:heis-matrix}
 \rho(c,a,b)e_x=\omega^{c+b\cdot x}e_{x+a},
 \qquad \omega=e^{2\pi i/p}.
\end{equation}
These $pN^2$ matrices are unitary and faithful.

\begin{theorem}[Odd extraspecial formula]\label{nd:central_extensions:thm:extra-odd}
For the module \eqref{nd:central_extensions:eq:heis-matrix},
\begin{equation}\label{nd:central_extensions:eq:extra-series}
 \boxed{\MGF_{E_n(p),V}
 =N^{-2}\PP_p\prod_i(1-t_i)^{-N}
 +(1-N^{-2})\prod_i(1-t_i^p)^{-N/p}.}
\end{equation}
Consequently,
\begin{align}\label{nd:central_extensions:eq:extra-coeff}
 [m_\tau]\MGF_{E_n(p),V}
 =\mathbf1_{p\mid|\tau|}\biggl[&N^{-2}
 \prod_i\binom{N+\tau_i-1}{\tau_i}\\
 &+(1-N^{-2})\mathbf1_{p\mid\tau_i\ \forall i}
 \prod_i\binom{N/p+\tau_i/p-1}{\tau_i/p}\biggr].\nonumber
\end{align}
If scalar $s$th roots are adjoined with $p\mid s$, then
\begin{equation}\label{nd:central_extensions:eq:central-product}
 \MGF_{\mu_sE_n(p),V}=\PP_s\MGF_{E_n(p),V}.
\end{equation}
\end{theorem}

\begin{proof}
The center consists of $p$ scalar matrices $\omega^cI$, a fraction $N^{-2}$
of the group.  Averaging their determinant factors gives the first term of
\eqref{nd:central_extensions:eq:extra-series}.  For a noncentral element, the weighted shift
\eqref{nd:central_extensions:eq:heis-matrix} has all $p$th roots of unity as eigenvalues, each with
multiplicity $N/p$.  Therefore
$\det(I-tg)=(1-t^p)^{N/p}$, giving the second term.  Extracting coefficients
with $(1-t)^{-a}=\sum_{d\geq0}\binom{a+d-1}{d}t^d$ proves
\eqref{nd:central_extensions:eq:extra-coeff}.  Equation \eqref{nd:central_extensions:eq:central-product} is
Theorem~\ref{nd:central_extensions:thm:projector}.
\end{proof}

The notebook constructs $E_1(3)$, $E_2(3)$, and $E_1(5)$, of respective
orders $27$, $243$, and $125$ and dimensions $3$, $9$, and $5$.  Every
coefficient through total degree six agrees with \eqref{nd:central_extensions:eq:extra-coeff}.
At $(t_1,t_2)=(0.07,0.11)$ the full determinant averages are:
\begin{center}
\begin{tabular}{@{}c c c c@{}}
\toprule
group & dense matrix average & formula & absolute error\\
\midrule
$E_1(3)$ & 1.006152412632117 & 1.006152412632117 & $7.8\times10^{-18}$\\
$E_2(3)$ & 1.016020200753344 & 1.016020200753347 & $2.7\times10^{-15}$\\
$E_1(5)$ & 1.000511444770786 & 1.000511444770785 & $1.6\times10^{-15}$\\
\bottomrule
\end{tabular}
\end{center}

The trace is zero off the center and equals $N\omega^c$ on the central
element $\omega^cI$.  Hence
\begin{equation}\label{nd:central_extensions:eq:trace-spike}
 \mathbb E[T^a\overline T^{\,b}]
 =\mathbf1_{p\mid(a-b)}N^{a+b-2},
\end{equation}
so $\mathbb E|T|^2=1$ while
$\mathbb E|T|^{2k}=N^{2k-2}$ for $k\geq2$.  Thus $T\to0$ in probability
as $n\to\infty$, but its higher balanced moments diverge: the limit is a
rare central spike, not a Gaussian.

For fixed $\tau$ and $d=|\tau|$, the central term of
\eqref{nd:central_extensions:eq:extra-coeff} is asymptotic to
$N^{d-2}/\prod_i\tau_i!$.  When all $\tau_i$ are divisible by $p$, the
noncentral term is asymptotic to
$(N/p)^{d/p}/\prod_i(\tau_i/p)!$.  Raw coefficients therefore generally
have no finite coefficientwise limit.

\subsection{Extraspecial \texorpdfstring{$2$}{2}-groups: separate plus and minus sectors}

Let $E_n^\varepsilon$ be extraspecial of order $2^{1+2n}$, sign
$\varepsilon\in\{+1,-1\}$, and let $N=2^n$.  Central elements again give
the root-filtered sector.  Noncentral elements whose squares are $+I$ have
eigenvalues $\pm1$ equally often; those whose squares are $-I$ have
eigenvalues $\pm i$ equally often.  Counting the two quadratic-form types
gives
\begin{equation}\label{nd:central_extensions:eq:extra-two}
 \boxed{\begin{aligned}
 \MGF_{E_n^\varepsilon}={}&
 2^{-2n}\PP_2\prod_i(1-t_i)^{-N}\\
 &+\left(\frac12+\varepsilon2^{-n-1}-2^{-2n}\right)
       \prod_i(1-t_i^2)^{-N/2}\\
 &+\left(\frac12-\varepsilon2^{-n-1}\right)
       \prod_i(1+t_i^2)^{-N/2}.
 \end{aligned}}
\end{equation}
The final two sectors can cancel in low degree; they may not be combined by
discarding the sign.

The rank-one plus group is $D_8$, represented by $\pm I,\pm X,\pm Z,\pm XZ$.
The rank-one minus group is $Q_8$, represented by
$\pm I,\pm iX,\pm iZ,\pm(-XZ)$.  Both $8$-element sets pass exact closure.
All coefficients through degree six agree with \eqref{nd:central_extensions:eq:extra-two}.  At the
same determinant test point, the errors are below $3.6\times10^{-18}$.

\subsection{Clifford groups and determinant-one lifts}

The full unitary Clifford normalizer contains $U(1)I$, so
Corollary~\ref{nd:central_extensions:cor:u1} gives
\[
 \MGF_{\mathrm{Cliff}_n}=1
\]
for the one-sided series.  A finite phase lift with scalar image $\mu_s$
instead applies $\PP_s$ and has lift-dependent coefficients in degrees
divisible by $s$.

Let $K_n=\mathrm{Cliff}_n\cap SU(N)$ with $N=2^n$.  It contains
$\{\zeta I:\zeta^N=1\}\cong\mu_N$, and therefore
\begin{equation}\label{nd:central_extensions:eq:cliff-selection}
 [m_\tau]\MGF_{K_n}=0\qquad\text{unless }N\mid|\tau|.
\end{equation}
For every fixed positive degree, eventually $N>|\tau|$, whence
\begin{equation}\label{nd:central_extensions:eq:cliff-limit}
 \MGF_{K_n}\longrightarrow1\qquad\text{coefficientwise}.
\end{equation}

In the one-qubit case, $K_1$ is the binary octahedral group already tested.
Adjoining $\mu_4$ produces $96$ distinct unitary matrices; all coefficients
through total degree six equal the binary-octahedral coefficient multiplied
by $\mathbf1_{4\mid|\tau|}$.  This directly tests the finite-lift rule.

Balanced quantities are the same for the full normalizer, any finite scalar
lift, $K_n$, and the projective quotient.  Multiqubit Clifford groups are
unitary $3$-designs, so their balanced series agrees with Haar unitary
through balanced degree three and differs at degree four.  The one-qubit
row in the table above gives the concrete values $1,2,5,15$ versus Haar
$1,2,5,14$.

More generally, suppose $\delta:\widetilde G\to\mu_q$ is a character and
$K=\ker\delta$.  Define
\[
 \MGF_{\widetilde G}^{(j)}
 =\frac1{|\widetilde G|}\sum_{g\in\widetilde G}\delta(g)^j
 \prod_i\det(I-t_i\rho(g))^{-1}.
\]
Character orthogonality gives the exact determinant-one formula
\begin{equation}\label{nd:central_extensions:eq:kernel-fourier}
 \boxed{\MGF_{K,V}=\sum_{j=0}^{q-1}\MGF_{\widetilde G}^{(j)}.}
\end{equation}
Indeed, $\sum_j\delta(g)^j$ equals $q$ on $K$ and zero off $K$, while
$|\widetilde G|=q|K|$ when $\delta$ is onto.  Passing to determinant one
therefore adds determinant-twisted sectors; it does not merely ``remove
scalars.''

\subsection{Sporadic covers: exact data workflow}

For a genuine irreducible module of a sporadic cover $2.G$, $3.G$, or $6.G$,
let $r$ be the order of its central scalar character.  Then
\[
 [m_\tau]\MGF_{\widetilde G,V}=0\qquad\text{unless }r\mid|\tau|,
\]
and the complete exact answer is \eqref{nd:central_extensions:eq:class-formula}.  An auditable
calculation imports three items from the represented covering group:
\begin{enumerate}
\item conjugacy-class sizes;
\item the genuine character values on those classes;
\item every class power map needed to the chosen truncation degree.
\end{enumerate}
The verification then compares a class-compressed sum against an element
enumeration when matrices are available, or against an independent character
table implementation otherwise.  There is no intrinsic rank asymptotic for
an isolated sporadic group.  Useful comparisons are the first separating
degree, balanced design strength, and equality of charge-zero series across
different lifts.

No external sporadic character table was supplied with this workspace, so
the notebook deliberately does not invent a sporadic numerical example.  The
binary octahedral and finite Weil cases exercise the same class/power and
phase mechanisms with explicit matrices; inserting certified ATLAS data in
\eqref{nd:central_extensions:eq:class-formula} is the correct sporadic verification strategy.

\subsection{Verification method and reproducibility}

The companion files are
\begin{itemize}
\item \texttt{verification.py}: constructions and comparison routes;
\item \path{marimo_notebooks/central_extensions.py}: the marimo interface,
grouped by group family;
\item \texttt{test\_verification.py}: a headless regression test.
\end{itemize}

For each dense matrix $g$, the code computes
$\Tr(\Sym^d g)$ from the recurrence
\begin{equation}\label{nd:central_extensions:eq:newton}
 h_0(g)=1,\qquad
 d\,h_d(g)=\sum_{k=1}^{d}\Tr(g^k)h_{d-k}(g).
\end{equation}
Thus
\[
 [m_\tau]\MGF=\frac1{|G|}\sum_{g\in G}\prod_i h_{\tau_i}(g)
\]
is obtained directly from actual matrix powers.  It is compared with
\eqref{nd:central_extensions:eq:extra-coeff}, coefficient extraction from \eqref{nd:central_extensions:eq:extra-two}, or
the scalar projector.  This is independent of the determinant-sector
comparison, which evaluates the full rational functions at
$(0.07,0.11)$.

For the projective tests, frame potentials are computed from the generated
matrices while the Weil comparator counts finite symplectic orbits.  The
binary octahedral matrices are also phase-twisted element by element.  Matrix
audits check expected order, distinctness, identity, unitarity, and exact
closure where a quadratic scan is modest.  The $243$-element $E_2(3)$ case
uses the exact Heisenberg multiplication law from \eqref{nd:central_extensions:eq:heis-matrix}
instead of a redundant $243^2$ dense closure scan.

\begin{center}
\begin{tabularx}{\linewidth}{@{}l r r r X@{}}
\toprule
family & groups & dimensions & checks & independent comparator\\
\midrule
odd extraspecial & 3 & $3,9,5$ & $90+3$ & spectral sectors and binomial formula\\
extraspecial $2$ & 2 & $2,2$ & $60+2$ & signed $\pm1/\pm i$ sectors\\
Schur/scalar lift & 2 & $2,2$ & 30 & homogeneous $\mu_4$ projector\\
projective spin/Clifford & 1 image & 2 & 5 & phase changes and Haar values\\
finite Weil & 2 & $3,5$ & 4 & symplectic orbit enumeration\\
\midrule
total & 9 constructions & $2$--$9$ & 194 & all pass\\
\bottomrule
\end{tabularx}
\end{center}

Here ``194'' means 180 coefficient, five determinant, and nine balanced
comparisons.  The maximum residual, $6.28\times10^{-12}$, comes from
floating-point roots of unity in an integer-valued coefficient; all stated
integer answers are within the declared $2\times10^{-9}$ tolerance.

\subsection{Asymptotic taxonomy and conclusions}

\begin{center}
\small
\begin{tabularx}{\linewidth}{@{}l X X@{}}
\toprule
family & one-sided behavior & balanced/projective behavior\\
\midrule
$2.S_n,2.A_n$ spin & odd degrees vanish & basic-spin fourth moment is linear in $n$\\
spin Weyl covers & parity projection & exterior-algebra centralizer law\\
sporadic covers & degree modulo central character & exact finite class formula; no intrinsic tower\\
finite Weil modules & lift-dependent scalar rule & fixed-$k$ symplectic orbit law stabilizes\\
extraspecial $p$-groups & coefficients generally diverge & rare central spikes; high moments diverge\\
full Clifford normalizer & exactly $1$ & Haar agreement through design degree\\
determinant-one Clifford & coefficientwise limit $1$ & same law as the projective quotient\\
scalar quotient & not canonically defined & charge zero is canonical\\
\bottomrule
\end{tabularx}
\end{center}

Central extensions therefore produce three distinct mechanisms:
\[
\boxed{
\begin{array}{ll}
\textbf{finite scalar center:}&\text{congruence projection on total degree};\\[1mm]
\textbf{growing or continuous scalar center:}&\text{one-sided collapse to }1;\\[1mm]
\textbf{projective representation:}&\text{information survives canonically at charge zero.}
\end{array}}
\]
The natural asymptotic objects are consequently projective or charge-graded
$\Lambda$-distributions, not phase-dependent one-sided distributions on a
scalar quotient.
